% ============================================================================
% Fichier  : partie4/P04_C11_controles_application.tex
% Titre    : Contrôles d'Application de la Dépense et de la Recette
% Auteur   : MINKA MI NGUIDJOÏ Thierry Emmanuel
% Version  : 0.1 -- novembre 2026
% ============================================================================
% Sources : INF3226 C18 (applications et données, resserré) ; GUID 5100
% (gestion des contrôles d'applications) ; COSO principe 10.
% Faits établis ici (LAB-09) :
%  MINÉQUIR, 9 611 mandats 2025 : liquidation > engagement : 0 ;
%   212 mandats > 50 000 000 FCFA, dont 4 sans visa du contrôle financier ;
%   doublons même fournisseur + même montant : 3 paires (P00) ; avec le
%   numéro de facture normalisé en plus : 1 paire (BATIRURAL) ; avec le
%   numéro brut : 0 (c'est pourquoi le contrôle de SIDÉPENSE n'a rien vu).
%  e-RECETTES 2025 : quittances Q-000001 à Q-018640 ; 37 numéros absents ;
%   112 annulations, dont 61 par le même agent (FOMEKONG).
% ============================================================================

\chapter{Contrôles d'Application de la Dépense et de la Recette}
\label{chap:applications}
\chaptermark{Contrôles d'application}

\epigraph{%
    « L'application fait ce qu'on lui a dit~; l'auditeur vérifie ce qu'on
    lui a dit. »%
}{Principe du cours}

\niveauChapitre{Intermédiaire}{Chapitres~\ref{chap:eval-sci}
et~\ref{chap:acces}.}

\begin{boxobjectifs}
\begin{itemize}
    \item Classer les contrôles d'application~: saisie, traitement,
          interface, sortie.
    \item Identifier les contrôles d'application attendus sur la chaîne de
          la dépense et sur celle de la recette.
    \item Tester un contrôle automatique, en sachant à quelle condition un
          essai unique suffit.
    \item Tester un contrôle par l'analyse de toute la population quand
          l'essai n'est pas possible.
    \item Coter les contrôles d'application et en tirer une conclusion.
\end{itemize}
\end{boxobjectifs}

\bigskip

% ============================================================================
\section*{Pourquoi ce chapitre}
% ============================================================================

Une fois les contrôles généraux évalués, l'auditeur sait s'il peut se
fier aux contrôles intégrés à l'application. Ce chapitre les passe en
revue sur la dépense (SIDÉPENSE) et sur la recette (e-RECETTES), et
montre comment les tester sans jamais toucher à la production.

\separateur

% ============================================================================
\section{Quatre familles de contrôles d'application}
\label{sec:c11-familles}
% ============================================================================

\begin{tableaucours}{L{2.4cm}L{4.6cm}L{5.2cm}}{%
    \tableauHeader{Famille} &
    \tableauHeader{Objectif} &
    \tableauHeader{Exemples}
}
Saisie & Seules des données valides entrent & Champs obligatoires,
formats, listes de valeurs, doublons, plafonds \\
Traitement & Les calculs et les enchaînements sont justes & Montant
liquidé au plus égal à l'engagement, calcul des pénalités \\
Interface & Les données transmises arrivent toutes, et intactes &
Totaux de contrôle entre SIDÉPENSE et le poste comptable \\
Sortie & Les états sont complets et ne vont qu'aux destinataires prévus
& Numérotation continue des quittances, états signés \\
\end{tableaucours}
\vspace{-0.8em}
\begin{center}{\small\itshape Tableau~11.1. Les quatre familles de contrôles
d'application}\end{center}

\separateur

% ============================================================================
\section{Sur la chaîne de la dépense}
\label{sec:c11-depense}
% ============================================================================

\begin{tableaucours}{L{4.0cm}L{2.0cm}L{6.2cm}}{%
    \tableauHeader{Contrôle attendu} &
    \tableauHeader{Famille} &
    \tableauHeader{Ce qu'il empêche}
}
Engagement au plus égal au crédit disponible & Saisie & Le dépassement
de crédit \\
Liquidation au plus égale à l'engagement & Traitement & Le paiement d'un
montant non engagé \\
Blocage des doublons (fournisseur, montant, facture normalisée) & Saisie
& Le double paiement \\
Visa du contrôle financier au-delà d'un seuil & Traitement & Le paiement
d'un mandat important sans visa \\
Imputation cohérente avec la nature de la dépense & Saisie & Les
imputations de complaisance \\
Totaux de contrôle de la transmission au poste comptable & Interface &
La perte ou l'ajout d'un mandat en cours de route \\
\end{tableaucours}
\vspace{-0.8em}
\begin{center}{\small\itshape Tableau~11.2. Les contrôles d'application de
la dépense}\end{center}

\begin{boxdef}
\termecle{Normalisation}\index{normalisation}~: transformation d'une
donnée avant comparaison (suppression des espaces, des tirets, des
suffixes comme «~bis~», passage en majuscules), pour que deux valeurs
qui désignent la même chose soient reconnues comme égales. Un contrôle de
doublon sans normalisation se contourne par une lettre.
\end{boxdef}

\separateur

% ============================================================================
\section{Sur la chaîne de la recette}
\label{sec:c11-recette}
% ============================================================================

\begin{tableaucours}{L{4.4cm}L{2.0cm}L{5.8cm}}{%
    \tableauHeader{Contrôle attendu} &
    \tableauHeader{Famille} &
    \tableauHeader{Ce qu'il empêche}
}
Numérotation continue et automatique des quittances & Sortie & Une
recette encaissée sans quittance enregistrée \\
Annulation d'une quittance sur motif et validation d'un second agent &
Traitement & L'annulation de complaisance \\
Rapprochement quotidien des encaissements et des versements & Interface &
Une recette non reversée \\
Calcul automatique des pénalités selon le taux en vigueur & Traitement &
Les remises non autorisées \\
\end{tableaucours}
\vspace{-0.8em}
\begin{center}{\small\itshape Tableau~11.3. Les contrôles d'application de
la recette}\end{center}

\begin{boxscenario}{Dossier e-RECETTES. Trous et annulations}
\textbf{Constatation.} En 2025, les quittances vont de Q-000001 à
Q-018640~; 37 numéros n'existent pas dans la base. Sur 112 annulations,
61 ont été faites par le même agent, M.~FOMEKONG, sans validation d'un
second agent~: e-RECETTES ne l'exige pas.

\medskip
\textbf{Analyse} (confiance moyenne). Les 37 numéros absents peuvent
correspondre à des quittances émises puis supprimées, ce que la
numérotation continue devait empêcher. La concentration des annulations
sur un agent n'est pas une preuve de fraude, mais un indice qui justifie
l'examen pièce par pièce des 61 annulations
(Chapitre~\ref{chap:donnees-fraude}).
\end{boxscenario}

\separateur

% ============================================================================
\section{Tester un contrôle automatique}
\label{sec:c11-tester}
% ============================================================================

\subsection{L'essai, si les contrôles généraux tiennent}

Un contrôle automatique fait la même chose à chaque opération. Si les
contrôles généraux sont efficaces (le programme et le paramètre n'ont pas
changé pendant la période), un essai qui doit passer et un essai qui doit
être bloqué suffisent, en environnement de test, avec un paramétrage
identique à la production.

\subsection{L'analyse de la population, sinon}

Si les contrôles généraux ne tiennent pas, ou s'il n'existe pas
d'environnement de test, on teste le contrôle par son \emph{résultat}~:
on cherche, dans toute la population, les opérations que le contrôle
aurait dû bloquer. Zéro opération trouvée sur toute l'année est une preuve
forte de fonctionnement~; une seule suffit à coter \NON{}.

\begin{boxoutils}{Tester deux contrôles de SIDÉPENSE par la population}
Feuille des 9\,611 mandats~: D montant liquidé, E montant engagé, F
montant, G visa (\texttt{OUI} ou vide).
\begin{lstlisting}
Liquidation > engagement :
H2 : =SI(D2>E2;"ECART";"")
Mandats > 50 M sans visa :
I2 : =SI(ET(F2>50000000;G2<>"OUI");"SANS VISA";"")
\end{lstlisting}
\end{boxoutils}

\begin{boxscenario}{Dossier MINÉQUIR. Deux contrôles par leur résultat}
\textbf{Constatation.} Aucun des 9\,611 mandats n'a une liquidation
supérieure à l'engagement. Sur les 212 mandats de plus de
50\,000\,000~FCFA, 4 n'ont pas de visa du contrôle financier.

\medskip
\textbf{Analyse} (confiance moyenne). Le contrôle de la liquidation a
fonctionné toute l'année (\OUI{}), mais la confiance reste moyenne~: les
accès n'étant pas maîtrisés, un compte ADM aurait pu le désactiver puis
le réactiver. Le contrôle du visa n'est pas bloquant (\NON{})~: les 4
mandats sont à examiner un par un.
\end{boxscenario}

\begin{boxobjet}{C-APP-03 · Blocage des doublons}
\textbf{Objectif.} Une même facture n'est payée qu'une fois.

\textbf{Critère.} SIDÉPENSE refuse un mandat de même fournisseur, même
montant et même facture normalisée qu'un mandat existant.

\textbf{Test de conception.} Le contrôle porte sur le numéro de facture
brut, sans normalisation.

\textbf{Test de fonctionnement.} Analyse des 9\,611 mandats~: avec le
numéro brut, aucune paire~; avec le numéro normalisé, une paire, celle du
prologue (les deux autres paires de même fournisseur et même montant ont
des factures différentes, justifiées par des livraisons distinctes).

\textbf{Cote.} \NON{}.
\end{boxobjet}

\begin{boxlecture}{Contrôles d'application de SIDÉPENSE}
\textbf{Constatation.} Sur six contrôles attendus (Tableau~11.2), trois
fonctionnent~: crédit, liquidation, transmission~; trois non~: doublons,
visa, imputation (non bloquante).

\textbf{Analyse} (confiance moyenne). Les contrôles qui fonctionnent
restent fragiles tant que la déficience majeure des accès n'est pas
corrigée.

\medskip
\hfill\conformite{3}{6}{0}
\end{boxlecture}

\begin{boxmetier}{Vérificateur}
Avant tout essai, demandez s'il existe un environnement de test dont le
paramétrage est identique à la production. S'il n'en existe pas, testez
par la population~: c'est plus long, mais c'est sans risque pour
l'entité.
\end{boxmetier}

\begin{boxmetier}{Contrôleur financier}
Si votre visa n'est pas bloquant dans l'application, chaque mandat qui
passe sans lui est un mandat que vous n'avez pas contrôlé. Demandez que le
seuil soit paramétré comme contrôle bloquant.
\end{boxmetier}

\begin{boxmetier}{Juge des comptes}
Une numérotation de quittances avec des trous est un indice
d'exhaustivité défaillante des recettes. Demandez l'explication de chaque
numéro absent.
\end{boxmetier}

\separateur

\begin{boxresume}
\begin{itemize}
    \item \textbf{Quatre familles.} Saisie, traitement, interface, sortie.
    \item \textbf{Dépense~:} crédit, liquidation, doublons, visa,
          imputation, transmission.
    \item \textbf{Recette~:} numérotation continue, annulations validées,
          rapprochement, pénalités.
    \item \textbf{Normaliser avant de comparer}~: un doublon se cache
          derrière une lettre.
    \item \textbf{Un essai suffit} si les contrôles généraux tiennent~;
          sinon, tester par la population.
    \item \textbf{Zéro sur toute la population} est une preuve forte~; un
          seul cas suffit à coter \NON{}.
\end{itemize}
\end{boxresume}

\begin{boxexo}[title={\ding{45}\quad Exercice~11.1. Classer \niveaubase}]
Rattachez à une famille~: (a) champ RIB obligatoire~; (b) total des
mandats transmis égal au total reçu~; (c) calcul de la pénalité de
retard~; (d) quittance imprimée avec un numéro.
\end{boxexo}

\begin{boxexo}[title={\ding{45}\quad Exercice~11.2. Normaliser \niveaubase}]
Écrivez la formule de tableur qui transforme «~F-2025-117 bis~» et
«~F2025117~» en une même valeur.
\end{boxexo}

\begin{boxexo}[title={\ding{45}\quad Exercice~11.3. Les 37 numéros \niveaumoyen}]
Décrivez votre démarche pour les 37 numéros de quittance absents
d'e-RECETTES~: pièces demandées, personnes interrogées, et ce qui ferait
passer l'indice au constat.
\end{boxexo}

\begin{boxexo}[title={\ding{45}\quad Exercice~11.4. Essai ou population \niveauexpert}]
La DAF vous propose de tester le nouveau blocage des doublons par un essai
en production, «~avec un mandat fictif qu'on annulera~». Dites pourquoi
vous refusez, ce que vous proposez, et à quelles conditions vous
conclurez \OUI{} pour l'année 2026.
\end{boxexo}

\begin{boxinfo}
\textbf{Transition.} L'analyse de toute la population a montré sa
puissance. Le Chapitre~\ref{chap:donnees-fraude} en fait un outil à part
entière, pour la paie et les marchés, et pour la détection de la fraude.
\end{boxinfo}
